I have far more people to thank than these pages can reasonably accommodate. In fact, I could probably write another thesis just for that. But in the interest of keeping things readable, I will focus here on a few individuals and groups, while many others are acknowledged throughout without being named individually. If you are reading this, you will likely know where you belong. And if I read this again in fifty years, I may forget some names, but not the memories.

Finally, this is coming to an end. First, I would like to thank myself for being patient and perseverant enough to drag this across the finish line. To my parents, mom and dad, thank you for the constant reminder that you were fully prepared to fly over and make sure I finished this PhD one way or another. Your unwavering support, both financial and otherwise, made it possible for me to pursue this path while somehow maintaining a lifestyle that suggests I made better financial decisions than I actually did. You made it possible for me to pursue a PhD full-time, even though at this stage of life, a stable corporate career would have been the more expected path in our culture. And to my brother, thank you for providing just enough early-life motivation to keep me moving forward.

To my “SSE academic family”: Sara Rosengren, for being the kind of mentor who gently but persistently pushes things to completion. I still remember sitting with you going line by line, brand by brand, category by category through grocery retail data until things finally started to make sense. Rickard Sandberg, for being supportive, understanding, and for always saying yes, especially to my reimbursement requests. With you, it was R code line by line, from linear models to non-linear ones, until they finally converged. Together, you taught me how things work, and occasionally, what we could have done differently. And Emelie Fröberg, the first person I met before this journey even started, and someone who has been there ever since. From academia to credit risk and everything in between, thank you for being there throughout. I suspect this is not the end of that journey.

Hannes Datta, my academic “adopted” parent from Tilburg University, thank you for taking me in and teaching me to be (arguably) great in both academia and industry. I know you would probably tell me not to be sentimental, but I will make a small exception here. My time in academia had its lowest points, including moments where statistically significant effect sizes I found approached $e^{-6}$, and even after rounding still showed up as $0.0000^*$, but also some of its best, many of which came from working with you. Thank you for believing in me when I did not, for giving me the opportunity to be part of your team, and for treating supervision more like coaching. We may not have a published paper together, but I hope I am, in some way, a reflection of your mentorship. You taught me how to think academically while being ready for industry, which I found more valuable than any 4-star job market paper. I am quite certain the lifetime returns from this mentorship will exceed $e^6$  EUR.

I would like to thank a few academics I had the chance to learn from along the way, especially Rik Pieters, Els Gijsbrechts, Marnik Dekimpe, Koen Pauwels, and Max Pachali, for offering some of the most expensive things a PhD student could ask for: free reviews of my papers, access to some of the best quantitative courses, and honest advice on navigating PhD life. I am also especially grateful to Anne Ter Braak and Lily (Xuehui) Gao for being my mock opponents and for providing one of the best mock defenses I could have experienced. I also appreciate the opportunity to be part of the IJRM newsletter, working together with Renana Peres and Martin Schreier, which gave me the chance to have one-on-one conversations with many inspiring academics, more than I can reasonably list here, but all featured in the newsletter for those curious.

%including Julian Wichmann, Valeria Stourm, Wiebke Keller, Danit Ein-Gar, Ali Umut Güler, Alex Bleier, Beth Fossen, Michal Shapira, Sarah Gelper, Simon Blanchard, and Klaus Miller.

Thank you to everyone who has been part of my life along the way, even if only for a chapter. From high school, to my undergraduate years, to my time in Japan, to my master’s at Cornell, to my early work at the Fiscal Policy Office, and my research experience at IESE before arriving at SSE, each of these places, and the people in them, shaped me in ways that made me more resilient. I would also like to thank my housemates during those years, especially those in Tokyo, Ithaca, and Barcelona, as well as the friends I met in the corridor at Forum in Stockholm, many of whom have become some of the best people I know and remain close to this day, for putting up with me more than they probably signed up for. At the time, I thought things could not get much worse after going through all of that with them. As it turns out, I just ended up with more stories to share when things did. Next time we meet, please remember to call me doctor (with apologies to my medical doctor friends, who kept me properly diagnosed and saved me from Googling the many strange symptoms that seemed to develop while I was doing my Doctor of Philosophy).

Thank you to all the friends I met during my PhD at SSE, who shared the many kinds of fights along the way. We fought for better causes, trying to find a significant effect and fill a research gap. We fought for our own causes, mostly PhD wellbeing, going out for drinks and sharing both joy and pain. And at times, we fought our own demons, hoping to tame them, or at least keep them from getting out of control. We were not exactly well-funded, but we somehow kept buying each other drinks. We spent hours talking about research ideas and, more often than not, doing very little about them. We vented about life, also without always solving much, but that was part of how we got through it. Our group chats, starting from “PhD baby” to “survivors,” “war veterans,” and “doctor bro,” say more than enough about how things went. I am glad to have gone through all of this with you. And to some of you, my office husbands, thank you for making everyday life a bit more enjoyable. I appreciated the lunch dates, the beer dates, and occasionally being a very happy third wheel.

Thank you to the faculty, DMS, and the centers, Center for Retailing (CFR) and Center for Data Analytics (CDA), for letting me be part of this affiliation. My research did not always fit perfectly, and feedback was sometimes limited, but I still valued the theoretical discussions, despite the data limitations in my empirical work. I enjoyed working as a TA with Magnus Söderlund, Daniel Tolstoy, and Kaisa Koskela-Huotari, particularly during the Canvas $1.0$ era in COVID times, when things were chaotic and sometimes the mic did not work. TA work made my PhD life slightly more bearable financially. I would also like to thank Karl Strelis, Alexandra Svahn Alberts, and Anna Burell for their help with all the paperwork related to CFR and CDA. Elena Braccia, our PhD administrator at the time, for always being responsive and able to handle all my requests without delay. Elizabeth Briggs, for patiently answering my endless questions and helping with all document and scholarship-related matters. I owe you all a lot.

% And Agneta Carlin, for supporting all administrative requests along the way. 

To all the academic colleagues I met along the way, including those who have since left academia, whether during my visit to Tilburg or at conferences in Budapest, Nashville, Hamburg, Odense, Bucharest, and Antwerp, it has been a pleasure crossing paths with you. Thank you for the shared dinners, drinks, and the many moments of venting that made the process easier to get through. At this point, I am convinced that conferences are simply a bundled version of therapy, and far more effective. 

I cannot talk about this journey without mentioning my time as a research visitor in “Tilly.” I really appreciated how generous the marketing department there was, from dinners to free drinks and bitterballen whenever we had guest speakers, and the PhD camps that made the experience even more memorable. Evenings at places like The Cat’s Back were some of my favorite times with friends there. I also enjoyed long walks in the forest behind campus, often with a cigarette and a colleague, mostly talking about how to deal with endogeneity. Thank you as well to those who helped carry me back (literally) when I overestimated my limits. Some stories are probably best left undocumented.
% and to those who showed up during my presentations, even when the topic was far from your own area of research.

To my colleagues at Nordea, the first place I joined after my PhD and my first role in industry abroad, thank you for giving me the opportunity to experience the best of both worlds. It was a place where I could contribute conceptually and be heard, without always needing to frame everything as a research gap or a theoretical contribution. I am especially grateful to Rene Eriksen for taking me in and for valuing my skills and knowledge, and for the support and understanding while I was still completing my PhD. And to all my colleagues at Nordea, thank you for making my first step into industry such a positive experience. I have not regretted leaving academia. This experience showed me that my PhD is, in fact, useful beyond its own field.

To all my Thai friends, and aunties, in Stockholm, as well as my Thai friends in the Netherlands, thank you for being the light and warmth in these cold countries (where warm food is not always easy to find). Some of you I knew before arriving here, and some I was lucky to meet along the way, but together you make Stockholm and the Netherlands feel more like home than they probably should. We shared summer trips simply because we were all here, more or less in the same situation. We keep our Thai lifestyle alive by regularly going out to eat together, and complaining about where all our money goes. To those who crossed oceans to visit me, thank you for bringing the equator a little closer, and for bringing spices with you. And to my very few Swedish friends outside of the PhD circle, thank you for being there throughout and for including me. I truly appreciate our friendship, and I promise to return the favor with homemade Thai food and a future invitation to Bangkok.

I would also like to thank those who helped bring this thesis across the final stretch.
%, where everything came together in less than two weeks. 
In particular, Björn Ericsson for coordinating the printing process, Riety for helping design the cover, and professional manuscript editor Craig MacDonald, who supported the review of this thesis by editing the final copy for clarity and conformity to the university’s standards. Without this support, I would likely not have been able to defend on schedule.

Most importantly, this PhD would not have been possible without the foundations that value science and are willing to support doctoral education. I am sincerely grateful to Torsten Söderbergs Stiftelse for funding my four years of PhD on retail analytics. I would also like to thank Hans Werthénfonden and Ann-Margret och Bengt Fabian Svartz Stiftelse for sponsoring my research visit to Tilburg and the invaluable learning I brought back from it. I am grateful to Iwar Sjögrens stipendiefond, Stig Svensson minne, ICA-köpmannens fond, Stiftelsen Louis Fraenckels Stipendiefond, and Rune Höglund fond for supporting many of my travels for study and conferences. Finally, I would also like to thank CFR and CDA for supporting my final year. I hope to contribute back in one way or another, or at the very least through taxes.

Dominik Andraszek, my love, my friend, my therapist, my doctor, my butler, my partner. Thank you for being everything and for going through all my ups and downs with me. It is not always easy to be with me, I sometimes cannot even stand myself, but you can, and I truly appreciate that. You have been there throughout, supporting me in ways I probably do not say often enough. I still struggle with the word “love” because, as an academic, I feel the need to define it before using it. But if love means someone I would like to spend my life with, someone I feel completely comfortable with, and someone I can share my daily life with 24/7 for more than 14 days without things falling apart, or at least without me causing too much trouble, then that is you. Without you, this PhD would probably have taken ten years, or more.